\begin{lemma}[Minimum bad pair: no repeated color on the path suffix]
Let \((\varphi,T)\) be a minimum bad pair with \(\rho(T)>2\) and
\(\rho(T)<|E(T)|\).  It is also impossible that no color missing at the
final vertex \(y_n\) is missing at any \(y_i\) with
\(\rho(T)\le i<n\).
\end{lemma}
\begin{proof}
Put \(q=\rho(T)\).  By
\noderef{TashTreeEarlyBetaNormalization}, replace \(\varphi\), if
necessary, by a coloring \(\varphi'\) for which \((\varphi',T)\) is
still a minimum bad pair satisfying the no-repetition hypothesis, and
choose \(\beta\) and \(i\le q-2\) such that \(\beta\) is missing at
both \(y_i\) and \(y_n\).

Apply \noderef{TashTreeGammaEpsilonSetup} to this normalized pair.  It
gives colors \(\gamma,\varepsilon\) and an index \(m\le q-2\) such
that \(\gamma\ne\beta\), \(\varepsilon\ne\gamma\), \(\gamma\) is
missing at \(y_q\), \(\varepsilon\) is missing at \(y_m\), and both
colors are unused through \(y_q\).

If \(\varepsilon\) is missing at \(y_n\), invoke
\noderef{TashTreeControlledRepeatedSuffixContradiction} with
\(\alpha=\varepsilon\), \(\beta=\gamma\), \(i=m\), and \(j=q\).
All its hypotheses are exactly the facts just obtained, and it gives a
contradiction.  Hence \(\varepsilon\) is not missing at \(y_n\).
Since \(\beta\) is missing there, the two colors are distinct.

Let \(Q\) be the \((\varepsilon,\beta)\)-component rooted at \(y_n\).
If \(Q\) meets a displayed vertex of \(T_{q-1}\), then
\noderef{TashTreeNoRepeatedIntersectionShortening}, using the early
missing witnesses \(y_i,y_m\), produces a smaller bad pair and hence a
contradiction.  Otherwise \(Q\) avoids every displayed vertex of that
prefix.  The hypotheses assembled above then match
\noderef{TashTreeNoRepeatedDisjointReduction}, which swaps this disjoint
component and reduces to the controlled repeated-suffix contradiction.
This is again impossible.  The two alternatives exhaust the component,
so the no-repetition case cannot occur.
\end{proof}
